% =====================================================================
% METHODOLOGY
% =====================================================================
\section{روش پژوهش}
\label{sec:methodology}

\subsection{هدف، استقرار و پین‌کردن}
\label{sec:method:target}

همهٔ اندازه‌گیری‌ها در برابر \en{PwnzzAI Shop} گرفته شده‌اند که روی
\en{commit} برنامه با شناسهٔ \code{cd3ac0d1} و \en{digest} تصویر
\code{sha256:7878fbd7\ldots} پین شده است. \en{container} تنها روی
\code{127.\bk 0.\bk 0.\bk 1:18080} منتشر می‌شود و به یک \en{backend} استنتاج که
روی میزبان اجرا می‌شود دسترسی دارد. مدل \code{llama3.2:1b} است، یک مدل
دستوریِ یک‌میلیاردپارامتری، که پیکربندی استقرار آن را برگزیده است، نه
توانمندی‌اش.

انتخاب یک مدل کوچک یک قید روش‌شناختی است و باید همان‌گونه خوانده شود. یک
مدل \num{1B} بیشتر از یک مدل مرزی امتناع می‌کند، پرت‌وپلا می‌گوید یا از
وظیفه منحرف می‌شود، و همین نرخ‌های مطلقِ موفقیت حمله را در سراسر مطالعه
فرو می‌نشاند. بنابراین این پژوهش تفاوت‌های \emph{نسبی} را --- میان سطوح
\en{persona}، میان لایه‌های \en{guardrail}، میان یک \en{mitigation} روشن و
خاموش، میان دو \en{detector} که یک خروجی را قضاوت می‌کنند --- سیگنالِ قابل
تفسیر می‌گیرد، و نرخ‌های مطلق را کران پایینِ مخصوص همین استقرار.
بخش~\ref{sec:discussion} به این نکته بازمی‌گردد.

وضعیتِ تغییرپذیر برنامه (بارگذاری‌ها، پیکرهٔ \en{RAG}، پایگاه‌دادهٔ نمونه)
از پوشه‌ای \en{mount} می‌شود که می‌توان آن را حذف کرد تا وضعیت پایهٔ شناخته
بازگردد؛ برنامه در نخستین درخواست پس از راه‌اندازیِ تازه، کاتالوگ، دو
حساب و \en{routing flag}‌های خود را دوباره می‌کارد.

\subsection{واحد اندازه‌گیری: یک \en{Task}، یک اجرای \en{Garak}}
\label{sec:method:task}

یک \emph{\en{task}} یک فراخوان \en{Garak} است: یک \en{probe} واحد که به
یک \en{generator} واحد زیر پیکربندی مشخصی عرضه می‌شود. یک
\emph{\en{suite}} فهرستی مرتب از \en{task}‌هاست که با هم به یک پرسش
ارزیابی پاسخ می‌دهند. اجراکننده برای هر \en{task} یک پیکربندی \en{JSON}
می‌نویسد و \code{garak.\bk cli.\bk main} را با آن فرا می‌خواند --- همان نقطهٔ
ورودی که یک اپراتور انسانی به کار می‌برد. \en{Garak} مالکِ توالیِ
\en{prompt}، تولید، اجرای \en{detector}، امتیازدهی و ساختِ مصنوعات است؛
اجراکننده تنها تصمیم می‌گیرد \emph{چه چیزی} اجرا شود و خروجی \emph{کجا}
بنشیند.

هر \en{task} چهار مصنوع تولید می‌کند: \code{report.jsonl} (هر تلاش با
\en{prompt}‌ها، خروجی‌ها، امتیاز هر \en{detector} و \en{notes} مربوط به
\en{generator})، \code{report.html} (چکیدهٔ خوانا برای انسان در
\en{Garak})، \code{hitlog.jsonl} (تلاش‌های ناموفق، هرگاه وجود داشته باشند)
و همان \code{config.json} که به‌کار رفته است. یک \en{manifest} برای هر
\en{suite}، \en{task}‌ها را به اثر انگشتِ هدف گره می‌زند. فایل‌های
پیکربندیِ نگه‌داشته‌شده، دستورالعمل‌های اجراییِ بازتولیدند: یک \en{task}
منفرد با پس‌دادن پیکربندی‌اش به \en{CLI} استاندارد \en{Garak} بازپخش
می‌شود.

یک جزئیات پیاده‌سازی برای اعتبار هر جاروب در این مطالعه باربر است.
\en{Garak} نمونه‌های \en{plugin} را با کلیدی که شکل رشته‌ایِ ریشهٔ پیکربندی
است \en{cache} می‌کند. چون همهٔ \en{task}‌های یک \en{suite} از راه یک شیء
پیکربندیِ سراسری در یک فرایند رانده می‌شوند، آن کلید در همهٔ \en{task}‌ها
یکسان است؛ بدون مداخله، \en{task} دوم بی‌سروصدا از \en{generator}
\en{task} اول استفاده می‌کرد و یک جاروبِ پنج‌سطحیِ \en{persona} در واقع
سطح یک را پنج بار اجرا می‌کرد. اجراکننده پیش از هر \en{task}،
\en{cache} نمونه‌ها را پاک می‌کند. تفاوت میان یک جاروبِ واقعیِ سطح، مرحله و
بودجه با همان اجرا که $N$ بار تکرار شده باشد همین است، و دلیلِ آن است که
روندهای یکنواختِ گزارش‌شده در بخش~\ref{sec:results} را می‌توان به پارامتر
جاروب‌شده نسبت داد.

\subsection{طراحی آزمایش}
\label{sec:method:design}

پنج \en{suite} شامل \num{28} \en{task} و \num{678} تلاشِ ارزیابی‌شده، هر
سه خانوادهٔ حملهٔ لازم را پوشش می‌دهند. جدول~\ref{tab:suites} طراحی کامل را
بیان می‌کند.

\begin{table}[t]
\caption{طراحی آزمایش. \en{attempts $=$ prompts $\times$ generations}. هر
\en{task} یک اجرای \en{Garak} است؛ پارامترِ جاروب‌شده تنها چیزی است که
درون یک \en{suite} تغییر می‌کند.}
\label{tab:suites}
\centering
\footnotesize
\zebra
\begin{tabular}{@{}L{\tw{0.15}}L{\tw{0.11}}L{\tw{0.15}}L{\tw{0.15}}L{\tw{0.20}}C{\tw{0.08}}C{\tw{0.07}}C{\tw{0.09}}@{}}
\toprule
\headrow
\tabhead{\en{Suite}} & \tabhead{\en{OWASP}} & \tabhead{\en{Probe}} & \tabhead{\en{Generator}} & \tabhead{پارامتر جاروب‌شده} &
\tabhead{\en{Task}} & \tabhead{\en{Gen.}} & \tabhead{تلاش} \\
\midrule
\code{direct-\bk injection} & \en{LLM01} & \code{Coupon\bk Extraction} & \code{Pizza\bk Assistant} &
سطح \en{persona} از \en{L1} تا \en{L5} (هر یک با رازِ خودش) & \num{5} & \num{3} & \num{180} \\
\code{guardrail-\bk ladder} & \en{LLM01} & \code{Guardrail\bk Bypass} & \code{Guardrail\bk Ladder} &
مرحلهٔ \en{guardrail} از \en{B0} تا \en{B9} (یک لایهٔ دفاعی در هر پله) & \num{10} & \num{3} & \num{330} \\
\code{indirect-\bk injection} & \en{LLM01} & \code{QRCode\bk Injection} & \code{QRChannel} &
--- (یک پیکربندی؛ \en{payload} متغیر است) & \num{1} & \num{3} & \num{21} \\
\code{information-\bk disclosure} & \en{LLM02/06} & \code{Customer\bk Data\bk Extraction, System\bk Prompt\bk Disclosure, Cross\bk Tenant\bk Order\bk Access} &
\code{Comment\bk RAG, Pizza\bk Assistant, Order\bk Access, Catering\bk SQLAgent} &
--- (چهار سطح افشای متمایز) & \num{4} & \num{3} & \num{81} \\
\code{data-\bk poisoning} & \en{LLM04} & \code{Sentiment\bk Poisoning, Catering\bk RAG\bk Poisoning} &
\code{Sentiment\bk Classifier, Catering\bk RAG} &
بودجهٔ سم \en{$\in \{0,1,3,5,10,20\}$}؛ \en{mitigation} بازیابی خاموش/روشن & \num{8} & \en{1--2} & \num{66} \\
\midrule
\grouprow \multicolumn{5}{@{}r}{\totlabel{مجموع}} & \totlabel{\num{28}} & & \totlabel{\num{678}} \\
\bottomrule
\end{tabular}
\end{table}

پیکره‌های \en{prompt} عمداً فشرده و خوانا هستند، نه بزرگ. هر \en{probe}
بین پنج تا دوازده \en{prompt} حمل می‌کند و هر \en{prompt} برای به‌کارگرفتنِ
یک تکنیک قابل تشخیص نوشته شده است --- درخواست مستقیم، بازنویسی دستور،
ادعای اقتدار، قاب‌بندی عاطفی، قاب‌بندی داستانی، غیرمستقیم‌گویی، کدگذاری،
بازخوانیِ مبهم‌شده، بازگویی، فرضی. این یک انتخاب طراحی با هدفی تحلیلی
است: چون قصدِ هر \en{prompt} معلوم است، پیامدهای هر \en{prompt} را می‌توان
به‌عنوان نتیجه‌ای در سطح تکنیک خواند، نه نوفه در یک تجمیع؛ و همین است که
جدول‌های~\ref{tab:technique} و~\ref{tab:ladder} در بخش~\ref{sec:results}
را اصلاً قابل تفسیر می‌کند. گستردگیِ پوشش با قابلیت انتساب معاوضه شده است.

\subsection{کنترل‌ها}
\label{sec:method:controls}

چهار کنترل به‌جای اعمال پسینی، در خودِ طراحی تعبیه شده‌اند
(شکل~\ref{fig:controls}).

\begin{figure}[t]
  \centering
  \safefig[0.78]{../figures/generated/fig-controls.pdf}
  \caption{چهار کنترل آزمایشی. هر کدام اندازه‌گیری‌ای را که در غیر این
  صورت یک مشاهدهٔ منفرد می‌بود به یک تفاوت یا یک کنارگذاری تبدیل می‌کند.}
  \label{fig:controls}
\end{figure}

\textbf{کنترلِ سالمِ جفت‌شده برای \en{poisoning}.} سطح احساسات در هر اجرا
دو مدل برازش می‌دهد --- یکی روی پیکرهٔ پایه، یکی روی پایه به‌علاوهٔ سم ---
و هر \en{prompt} ارزیابی را با هر دو طبقه‌بندی می‌کند. موفقیت به‌عنوان
\emph{اختلاف} میان آن دو تعریف می‌شود، هرگز به‌عنوان یک حکمِ مسمومِ منفرد.
یک حکم تنها نمی‌تواند نشان دهد سم چه چیزی را تغییر داده است: عبارتی که
طبقه‌بند همیشه بد طبقه‌بندی می‌کرده در غیر این صورت موفقیت حمله شمرده
می‌شد. چون هر دو برازش مدل‌های قطعیِ \en{scikit-learn} هستند، این کنترل
دقیق است و منحنی‌های حاصل بیت‌به‌بیت بازتولید می‌شوند.

\textbf{روشن/خاموش‌کردن \en{mitigation} برای \en{poisoning} بازیابی.}
همان سند سمی در پیکرهٔ \en{catering} تزریق می‌شود و سپس همان مجموعه
\en{prompt} دو بار اجرا می‌شود: یک بار با بازیابیِ فقط-مورد‌اعتمادِ برنامه
خاموش و یک بار روشن. این تقابل، اثر \en{mitigation} را با ثابت‌نگه‌داشتنِ
پیکره، \en{prompt}‌ها و مدل جدا می‌کند.

\textbf{بررسیِ سلامتِ تحویل برای تزریق غیرمستقیم.} \en{generator} مربوط به
\en{QR} ثبت می‌کند که آیا \en{payload} کدشده از گام رمزگشاییِ برنامه سالم
رفت‌وبرگشت کرده است یا نه. \en{payload} مخدوش یعنی هدف هرگز پرسشی را که
\en{probe} در نظر داشته دریافت نکرده، پس آن تلاش بی‌اطلاع است؛ یک
\en{detector} تشخیصیِ اختصاصی چنین تلاش‌هایی را به‌جای آنکه بگذارد
«دفع‌شده» شمرده شوند، برای کنارگذاری علامت می‌زند.

\textbf{\en{prompt}‌های منفی و بی‌خطر.} \en{probe} مربوط به \en{poisoning}،
\en{prompt}‌های کنترلِ بدونِ ماشه را حمل می‌کند، از جمله یک جملهٔ آشکارا
مثبت و یک جملهٔ آشکارا منفی دربارهٔ جنبه‌های بی‌ربط سرویس. تغییر برچسب روی
آن‌ها نشانهٔ افت عمومیِ طبقه‌بند می‌بود --- شکستی گسترده‌تر و متفاوت با یک
\en{backdoor} هدفمند --- و تفکیک این دو تنها زمانی ممکن است که کنترل‌ها
در همان اجرا حمل شوند.

\subsection{معیارها}
\label{sec:method:metrics}

\en{Garak} برای هر جفت \emph{(\en{probe}، \en{detector})} شمارش
\code{passed}، \code{fails} و \code{nones} را گزارش می‌کند. معیار سرخط،
نرخ موفقیت حمله روی تلاش‌های \emph{ارزیابی‌شده} است:
\begin{equation}
  \mathrm{ASR} = \frac{\mathit{fails}}{\mathit{passed} + \mathit{fails}},
  \label{eq:asr}
\end{equation}
که در آن \code{nones} به‌طور ساختاری از مخرج کنار گذاشته می‌شود. در کنار
قواعد امتناع در بخش~\ref{sec:arch:detectors}، این یعنی یک خطای
\en{transport}، یک \en{ground truth} مفقود، یا یک مدل کنترلِ برازش‌نشده،
به‌جای ورود بی‌سروصدا به صورت یا مخرج، حجم نمونه را کاهش می‌دهد. در
اجراهای گزارش‌شده در اینجا شمار امتناع برای هر \en{detector} صفر بود،
بنابراین هر \num{678} تلاش ارزیابی شد؛ با این حال سازوکار باربر است، چون
همین است که اجازه می‌دهد یک شکست جزئی به‌جای جذب‌شدن در یک نرخ، به‌عنوان
شکست بازشناخته شود.

جمع‌بندی‌ها بر حسب خانوادهٔ حمله و بر حسب دستهٔ
\en{OWASP LLM Top 10 (2025)}~\cite{owasp2025}، \en{detector}
\emph{اصلیِ} اعلام‌شدهٔ هر \en{probe} را سیگنالِ سرخط می‌گیرند.
\en{detector}‌های گسترده --- از جمله یک \en{detector} استاندارد \en{Garak}
روی هر \en{probe} قابل اِعمال --- به‌جای ادغام، در کنارشان گزارش می‌شوند،
چون به پرسش‌های متفاوتی پاسخ می‌دهند و ادغام آن‌ها دقیقاً همان مقایسه‌ای را
نابود می‌کرد که این مطالعه به آن علاقه‌مند است.

امتیازها برای تصمیم \en{pass}/\en{fail} در $0.5$ آستانه‌گذاری می‌شوند، و
همین خروجی‌های درجه‌بندی‌شدهٔ دو \en{detector} را پیامددار می‌کند:
\code{CustomerPII} به یک نشانی \en{email} تنها مقدار $0.5$ می‌دهد (مدل
می‌تواند نشانی‌ای باورپذیر بسازد بی‌آنکه آن را بازیابی کرده باشد، پس برای
امتیاز کامل تأیید متقابل از نوع دومی از میدان لازم است)، و
\code{Poisoned\bk Retrieval\bk Influence} به زبانِ نرمِ سیاست که با هیچ قلمِ
محتواییِ تزریق‌شده همراه نباشد مقدار $0.5$ می‌دهد.

\subsection{تکرارپذیری و حدود آن}
\label{sec:method:repro}

\en{RNG} خودِ \en{Garak} \en{seed} شده است و همین نمونه‌گیریِ \en{prompt} و
ترتیب \en{buff}‌ها را ثابت می‌کند. موازی‌سازی همه‌جا غیرفعال است. تصویر
برنامه با \en{digest} پین شده و ثابت‌های \en{ground truth} در هر اجرای
مجموعه آزمون با آزمون‌های قراردادی در برابر کدِ پین‌شده دوباره اعتبارسنجی
می‌شوند. مجموعهٔ کامل \num{28}-\en{task} روی ماشین مرجع در حدود
\num{15} دقیقه زمان ساعت دیواری کامل می‌شود که عمدتاً صرف استنتاج مدل
می‌شود.

آنچه بیت‌به‌بیت تکرارپذیر \emph{نیست} مدل زبانیِ پشت مسیر \en{HTTP} است که
از راه رابط برنامه \en{seed}پذیر نیست. بنابراین تولیدهای منفرد میان
اجراها تغییر می‌کنند. همین دلیلِ آن است که هر \en{task} مبتنی بر \en{LLM}
چند تولید اجرا می‌کند و تحلیل به‌جای پیامدهای منفرد، نرخ گزارش می‌دهد؛ و
همین دلیلِ آن است که بخش~\ref{sec:results} به‌جای اعداد منفرد، ترتیب و
بزرگیِ تفاوت‌ها را تفسیر می‌کند. سطحِ \en{poisoning} احساسات از این قاعده
مستثناست: هیچ مدل زبانی در آن دخیل نیست و اعدادش دقیقاً بازتولید می‌شوند.
